\chapter{Designing a Universe from Scratch}

\begin{chaptermapbox}
\textbf{In this chapter:}
\begin{compactchecklist}
\item The design brief: what a universe-builder with finite resources would prioritize
\item Short rules, identical parts, and conserved quantities as engineering choices
\item Discrete space and time: what physics suggests, what it rules out, and what remains open
\item Singularities as crashes, and the guard clauses nature seems to install
\end{compactchecklist}
\end{chaptermapbox}

Part I ended at a wall. Let's walk around to the other side of it and try a thought experiment.

You've been hired to build a universe. Not to simulate ours, but to design one from nothing. Your client has a few requirements. The universe has to run for billions of years without crashing. It has to support complex structures: stars, chemistry, eventually something like life. It has to be internally consistent, so that no part of it contradicts any other part. And, most important for this book, you have a fixed budget. Finite memory. Finite processing power. Whatever you build has to fit.

Finite resources are a rule of this thought experiment. We have no measurements of the hardware or physical laws of a hypothetical external simulator.

What would you do?

This chapter walks through the design decisions a sensible engineer would make. Then, at each step, it holds those decisions up against the universe we actually have. The match isn't perfect. It's close enough to be worth staring at.

\section{Engineered Reality}

\subsection{Principle One: Keep the Rules Short}

The first thing any experienced programmer learns is that complexity is expensive. Every special case is a place for bugs to hide. Every exception has to be checked, every frame, forever.

So you'd want your universe's rules to be short. A handful of equations, applied everywhere, with no exceptions for particular places or times. The complexity of the world would come from running those simple rules on many parts, not from the rules themselves.

That's exactly what physics has found. The core of the Standard Model of particle physics, together with Einstein's equation for gravity, can be written compactly enough to fit on a coffee mug, and physics departments sell such mugs. Those few lines, plus about two dozen measured numbers, account for essentially every experiment ever done on Earth.

Computer scientists have a formal way to talk about this. The Kolmogorov complexity of a piece of data is the length of the shortest program that produces it \cite{livitanyi2008}. A string of a million random digits has high complexity; you can't describe it more briefly than by listing it. A string of a million zeros has very low complexity; ``print zero a million times'' does the job. The observed universe is staggeringly detailed, but the rules that generate its behavior appear to be very short. In the language of compression, the universe is extremely compressible.

\begin{researchnotebox}
\textbf{Well supported:} The laws of physics as currently known are compact, and they apply uniformly across space and time to the precision we can test. \textbf{Exploratory:} Calling this ``compressibility'' is a reading, not a measurement. The initial conditions of the universe, which the laws don't specify, may carry much of the information. Chapter~9 returns to that.
\end{researchnotebox}

\subsection{Principle Two: Reuse the Same Parts}

Game developers learned long ago that you don't store a separate, unique model for every tree in a forest. You store one tree, and then you place many copies of it, each with a position and maybe a slight color tweak. Software engineers call a version of this the flyweight pattern: share everything you can, and store only what differs.

In our universe, every electron is identical to every other electron. Not similar. Identical. They have exactly the same mass, the same charge, the same spin, everywhere and at every time anyone has checked. The same is true of every up quark, every photon, every neutrino of a given type. Nature has a very small catalog of part types, and it uses them over and over.

John Wheeler once called Richard Feynman with a playful idea: maybe all electrons are identical because there's only one electron, zig-zagging forward and backward through time so that it shows up everywhere \cite{feynman1965}. Feynman didn't take it literally, but it helped him think about antimatter. From an engineering standpoint, the joke points at something real. A universe that stores one definition of ``electron'' and stamps it out \(10^{80}\) times is far cheaper than one that stores \(10^{80}\) slightly different electrons.

This identity has measurable consequences. Because identical particles can't be told apart, swapping two of them doesn't create a new state. That cuts the number of distinct configurations a system can be in by a factor of \(N!\), which for even modest \(N\) is a huge saving. Physicists knew about this factor long before quantum mechanics; without it, the classical formula for the entropy of a gas gives wrong answers. It's called the Gibbs paradox, and the resolution is that nature doesn't keep separate labels for identical parts.

\subsection{Principle Three: Protect Invariants}

Good software has invariants, properties that must always be true no matter what the code does. A bank's ledger must always balance. A file system must never lose track of a block. If an invariant breaks, you know there's a bug.

Physics has invariants too. Energy is conserved. So is momentum, angular momentum, and electric charge. These aren't arbitrary rules stacked on top of the physics. In 1918, Emmy Noether proved that each conservation law follows from a symmetry \cite{noether1918}. If the laws don't change over time, energy is conserved. If they don't change from place to place, momentum is conserved. If they don't care which direction you face, angular momentum is conserved.

That's a beautiful piece of design. You don't have to write a separate check for each conservation law. You make the rules symmetric, and the invariants come for free. The symmetry tells you which quantity to check when a calculation goes wrong.

\subsection{Principle Four: Keep Things Local}

The final principle is locality. In a well-designed simulation, each piece of the world should only need information from its immediate neighbors to update itself. That way you can split the work across many processors, each handling its own patch, and they only need to talk at the edges.

Most of physics is local in exactly this sense. Fields at a point change based on fields at nearby points. Influences spread outward at finite speed. In 1969, the German engineer Konrad Zuse, who built some of the first programmable computers, took this idea to its conclusion and proposed that the universe might be a cellular automaton, a grid of cells, each updating based on its neighbors \cite{zuse1969}. His book was called \emph{Rechnender Raum}, ``Calculating Space.''

Quantum mechanics complicates the locality story in a fascinating way, through entanglement. That's the subject of Chapter~8. For now, note that the speed of light enforces a strict kind of locality on cause and effect, and Chapter~5 reads that speed as a synchronization limit.

\section{Discretized Space-Time}

\subsection{The Problem with Infinity}

Back to the design brief. You have finite memory. How do you represent position?

A generic real-valued position cannot be stored as a finite string of digits. Some values have short exact symbolic descriptions, but that does not solve arbitrary numerical evolution. A simulator must specify how much precision its question requires, and chaos can make trajectory predictions sensitive to that choice.

One engineering choice is a finite lattice with discrete updates. Within a bounded domain it gives finitely many position labels, although other quantities may still require approximation. This is a choice for the model, not a deduction about spacetime.

A screen has pixels and a displayed animation has frames. The underlying model can still contain positions and events finer than either. Display resolution is not the same as the resolution of the simulated world.

\subsection{Nature's Discreteness, Where We Know It Exists}

Physics has already found discreteness in many places where classical physics assumed continuity. That's where the word ``quantum'' comes from.

Energy in an atom comes in discrete levels. An electron in a hydrogen atom can sit in the lowest level or the second level, but nowhere in between. Light comes in discrete packets called photons. Angular momentum comes in whole or half-whole multiples of a fundamental unit, Planck's constant divided by \(2\pi\). Electric charge comes in fixed lumps; nobody has ever seen a free particle with half an electron's charge.

Quantized measurement outcomes do not imply a finite list of all possible quantum states: amplitudes can still vary continuously in the theory. Nor do discrete atomic energy levels establish discrete space or time.

\subsection{The Planck Scale}

In 1899, Max Planck noticed that by combining three fundamental constants, the speed of light \(c\), the gravitational constant \(G\), and his own new constant \(h\), he could build natural units of length, time, and mass that didn't depend on any human convention. Using the reduced constant \(\hbar = h/2\pi\), the Planck length is
\[
\ell_P = \sqrt{\frac{\hbar G}{c^3}} \approx 1.6 \times 10^{-35}\ \text{m},
\]
and the Planck time, the time light takes to cross one Planck length, is about \(5.4 \times 10^{-44}\) seconds \cite{codata2018}.

These scales are where, as Chapter~2 described, quantum mechanics and gravity both become important at once. Many physicists suspect that below the Planck length, the usual notion of distance stops making sense. One rough argument goes like this. To measure a smaller distance, you need a probe with a shorter wavelength, which means higher energy. Push the energy high enough to resolve a Planck length, and you've packed so much energy into so small a region that it collapses into a black hole, hiding the very distance you were trying to see.

That sounds a lot like a resolution cap. Ask for more detail than the system supports, and the system refuses to answer.

\subsection{What Quantum Gravity Suggests}

Some approaches to quantum gravity make discreteness explicit. In loop quantum gravity, Carlo Rovelli and Lee Smolin showed in 1995 that areas and volumes come in discrete steps, like the energy levels of an atom \cite{rovellismolin1995}. A surface can have one unit of area, or two, or some other value from an allowed list, but not just any value.

This is subtler than a pixel grid, and the difference matters. A naive grid would pick out preferred directions. A checkerboard looks different along its rows than along its diagonals. If space were a simple cubic lattice, light traveling diagonally might move at a slightly different speed than light traveling along the grid lines. In loop quantum gravity, areas and volumes are discrete but there's no fixed background grid, so the theory can, in principle, avoid preferred directions.

\subsection{What Experiments Rule Out}

Experiments constrain specific ways that a discrete model might alter observable physics.

\begin{researchnotebox}
\textbf{Well supported (observation):} Fermi's observations of the 2009 gamma-ray burst GRB~090510 constrain particular models in which photon speed varies with energy \cite{abdo2009}. The inference depends on the assumed dispersion law and on what can be inferred about emission times at the source. It is not a general test excluding discrete spacetime.

\textbf{Other tests.} The Fermilab Holometer, a pair of large laser interferometers, looked for a predicted ``holographic noise'' in spacetime and found none at its sensitivity \cite{chou2016}. Silas Beane and colleagues have proposed that, if the universe were computed on a cubic lattice, the highest-energy cosmic rays might show a telltale cutoff and directional pattern \cite{beane2014}. No such lattice signature has been confirmed.

\textbf{Status:} Whether spacetime is fundamentally discrete is \textbf{exploratory}. Some dispersion models are tightly constrained; discreteness itself remains open.
\end{researchnotebox}

A candidate discrete theory must reproduce the symmetries and propagation speeds that observations constrain. A pixel metaphor does not specify such a theory.

\subsection{A Different Kind of Cap}

There's another way to bound information without pixelating space, and it may be the one nature actually uses.

In 1981, Jacob Bekenstein proposed an entropy bound for suitable systems of given size and energy \cite{bekenstein1981}. Its motivation comes from black-hole thermodynamics. Applying information bounds outside their intended regimes requires care, especially in cosmology.

That's a memory cap rather than a resolution cap. It says: you can describe this region however you like, continuously or discretely, but the bound limits distinguishable information under its physical assumptions; it does not supply an exact finite-bit description of every quantum amplitude. Chapter~8 explores this idea in depth, because it leads to the holographic principle, one of the strangest and most important ideas in modern physics.

For the design brief, an entropy bound suggests another question to ask besides whether space has pixels: how much distinguishable information can a specified physical system support? Chapter~8 examines the assumptions behind the answer.

\begin{quickpracticebox}
Match the claim to its limit. Atomic energy levels provide a tested example of quantization; they do not establish a spatial grid. Discrete areas in a quantum-gravity proposal are a theoretical prediction awaiting confirmation. An entropy bound constrains information under stated conditions; it is not a picture of pixels. Which of these would actually support the claim you want to make?
\end{quickpracticebox}

\section{Elimination of Singularities}

\subsection{Guard Clauses}

Every programmer has written a guard clause. Before you divide by \(x\), you check that \(x\) isn't zero. Before you take a square root, you check that the number isn't negative. Before you allocate memory, you check that you have some left. Guard clauses aren't glamorous. They're what keep software from crashing.

Chapter~2 described the singularities of general relativity as places where the equations divide by zero. A universe-builder would want to prevent those, or at least contain them. Looking at physics through that lens, you find several mechanisms that act like guard clauses, some well established and some still hypothetical.

\subsection{The Speed Limit as a Soft Cap}

Take velocity. In Newton's physics, there's no limit to how fast something can go. Push it harder, it goes faster, forever. A simulation built on Newton's laws could, given enough force, send a particle to infinite speed in finite time, and the numbers would blow up.

Relativity installs a cap. As an object with mass approaches the speed of light, the energy needed to accelerate it further grows without bound. You can get to 99 percent of light speed, then 99.9 percent, then 99.99 percent, but each step costs more than the last, and \(c\) itself is never reached. In software terms, it's a soft cap with an asymptote, not a hard \texttt{if} statement. That's elegant, because there's no sudden discontinuity at the limit for the code to stumble over.

\subsection{Why Atoms Don't Collapse}

Here's a guard clause that physics discovered the hard way.

Around 1911, physicists knew that atoms had a tiny positive nucleus with electrons somewhere outside it. But classical physics predicted disaster. An electron circling a nucleus is an accelerating charge, and accelerating charges radiate energy. The electron should lose energy, spiral inward, and crash into the nucleus. The calculation gives a lifetime of about \(10^{-11}\) seconds. Every atom in the universe should have collapsed almost instantly.

Quantum mechanics fixed this. The uncertainty principle says you can't confine an electron to a very small region without giving it a large spread of momentum. Squeeze it toward the nucleus, and its kinetic energy rises faster than its electrical energy falls. There's a minimum-energy state, the ground state, and the electron can't go lower. The collapse is prevented by the same rule that makes the quantum world fuzzy.

There's a further layer of protection for bulk matter, which Chapter~5 covers in detail. In the 1960s and 1970s, physicists proved that ordinary matter is stable, meaning a pile of atoms doesn't collapse into an ever-denser blob, only because electrons obey the Pauli exclusion principle \cite{liebthirring1975}. Without that principle, matter would implode.

\subsection{Cosmic Censorship as a Sandbox}

Now the big one: black holes. General relativity predicts a true singularity at a black hole's center. Physics doesn't prevent that singularity from forming in the theory. What it does seem to do is hide it.

Every black hole we know of is surrounded by an event horizon, a boundary from which nothing, not even light, can escape. Whatever goes wrong at the singularity stays inside. In 1969, Roger Penrose proposed that this is always the case, an idea called the cosmic censorship conjecture \cite{penrose1969}. Nature, he suggested, never shows a ``naked'' singularity to the outside universe.

To a programmer, this looks like sandboxing. If you can't prevent a piece of code from crashing, you can at least run it in an isolated container, so that the crash doesn't take down the rest of the system. The event horizon acts like that container's wall. Outside observers see a well-behaved object with a mass, a spin, and a charge. Whatever singular behavior happens inside never feeds back out.

\begin{mythvsfactbox}
\textbf{Myth:} Physics has proven that singularities can't exist.

\textbf{Fact:} General relativity predicts singularities, and the cosmic censorship conjecture, that they're always hidden behind horizons, is unproven in general. Mathematicians have found special, finely tuned situations in the equations where naked singularities can form. Most physicists expect a theory of quantum gravity to replace the singularity with something finite, but that expectation hasn't been tested.
\end{mythvsfactbox}

\subsection{A Density Cap?}

If quantum gravity does remove singularities, the natural guess is that density tops out somewhere near the Planck density, about \(5 \times 10^{96}\) kilograms per cubic meter. That's the mass of roughly \(10^{66}\) Suns squeezed into a single cubic meter. Unimaginably dense, but finite.

Some models in loop quantum cosmology go further. They suggest that when the density of a collapsing universe approaches the Planck scale, quantum effects make gravity repulsive, and the collapse turns into a bounce \cite{ashtekar2006}. Instead of a Big Bang singularity, there would have been a ``Big Bounce'' from a prior contracting phase. These ideas are \textbf{exploratory}, and no observation yet distinguishes them from other possibilities. But from a design perspective, they're exactly the kind of mechanism you'd install: a maximum value, enforced by the rules themselves, that turns a potential crash into a recoverable event.

\section{Reading the Design Brief Backwards}

Let's step back and look at what we've done. We started with a set of engineering principles any careful universe-builder might follow:

\begin{compactchecklist}
\item Keep the rules short and uniform.
\item Reuse identical parts.
\item Derive invariants from symmetry so they can't break.
\item Keep updates local so work can be divided.
\item Cap the information any region can hold.
\item Guard against infinities, or at least sandbox them.
\end{compactchecklist}

And at each step, physics had already found something that does the same job. Compact laws. Identical particles. Noether's theorem. Local fields and a finite speed of influence. The Bekenstein bound. Relativistic speed limits, quantum ground states, the Pauli exclusion principle, and event horizons.

It's tempting to read this as evidence of design. I'd caution against that. There's a selection effect at work. A universe without stable atoms couldn't produce physicists to ask the question. A universe whose laws were long and full of exceptions might be impossible to discover laws for. We can only look at the universe we're in, and that universe was always going to have features that let it last long enough to be looked at.

These comparisons identify useful similarities. They do not show that nature has selected the cheapest implementation, or that the assumed budget explains its laws.

One engineering strategy remains: defer a calculation until its result is needed. Quantum measurement invites that comparison, though the fit needs care.

\begin{keytakeawaybox}
The design brief offers useful comparisons, but quantization, proposed spacetime discreteness, and entropy bounds are different claims with different evidence.
\end{keytakeawaybox}
